\section*{Capítulo \ref{ch:g11:trigo} --- Trigonometría: la circunferencia unidad}

\begin{solution}{exo:g11:trigo:1}
A \omterm{def:g11:trigo:radian}{radianes}: $30^\circ = \frac{\pi}{6}$, $45^\circ = \frac{\pi}{4}$,
$120^\circ = \frac{2\pi}{3}$, $270^\circ = \frac{3\pi}{2}$.

A grados: $\frac{\pi}{5} = 36^\circ$, $\frac{5\pi}{6} = 150^\circ$,
$\frac{7\pi}{4} = 315^\circ$, $\frac{2\pi}{9} = 40^\circ$.
\end{solution}

\begin{solution}{exo:g11:trigo:2}
$\frac{2\pi}{3}$ y $-\frac{\pi}{4}$ ya están en $\intoc{-\pi}{\pi}$:
segundo cuadrante y cuarto cuadrante, respectivamente.

$\frac{17\pi}{6} - 2\pi = \frac{5\pi}{6}$: segundo cuadrante, cerca del
semieje $x$ negativo.

$-\frac{7\pi}{2} + 4\pi = \frac{\pi}{2}$: lo más alto de la
circunferencia, $(0, 1)$.
\end{solution}

\begin{solution}{exo:g11:trigo:3}
$\cos\frac{3\pi}{4} = \cos\left(\pi - \frac\pi4\right)
= -\cos\frac\pi4 = -\frac{\sqrt2}{2}$.

$\sin\frac{5\pi}{6} = \sin\left(\pi - \frac\pi6\right)
= \sin\frac\pi6 = \frac12$.

$\cos\left(-\frac\pi3\right) = \cos\frac\pi3 = \frac12$.

$\sin\frac{4\pi}{3} = \sin\left(\pi + \frac\pi3\right)
= -\sin\frac\pi3 = -\frac{\sqrt3}{2}$.

$\cos\frac{11\pi}{6} = \cos\left(-\frac\pi6 + 2\pi\right)
= \cos\frac\pi6 = \frac{\sqrt3}{2}$.
\end{solution}

\begin{solution}{exo:g11:trigo:4}
$\sin^2 t = 1 - \cos^2 t = 1 - \frac{25}{169} = \frac{144}{169}$, luego
$\sin t = \pm\frac{12}{13}$. En $\intoo{-\frac\pi2}{0}$ (cuarto
cuadrante) la ordenada es negativa: $\sin t = -\frac{12}{13}$.
\end{solution}

\begin{solution}{exo:g11:trigo:5}
Usando la \cref{prop:g11:trigo:associated}:
\[
A = (-\cos x) + \cos x + \cos x - \cos x = 0 ,
\]
\[
B = \sin x + (-\sin x) + (-\sin x) = -\sin x .
\]
\end{solution}

\begin{solution}{exo:g11:trigo:6}
$\cos t = \frac{\sqrt3}{2}$: una solución es $\frac\pi6$, luego
$t = \pm\frac\pi6$ (los desplazamientos de $2k\pi$ salen del \omterm{def:g10:numbers:interval}{intervalo}).

$\sin t = \frac12$: $t = \frac\pi6$ o $t = \pi - \frac\pi6 =
\frac{5\pi}{6}$.

$\cos t = -1$: el único punto de \omterm{def:g10:coordgeom:system}{abscisa} $-1$ es $M(\pi)$, luego
$t = \pi$.
\end{solution}

\begin{solution}{exo:g11:trigo:7}
$\sin t = -\frac{\sqrt3}{2}$: partiendo de $a = -\frac\pi3$, las familias
son $-\frac\pi3 + 2k\pi$ y $\pi + \frac\pi3 + 2k\pi$; en
$\intco{0}{2\pi}$: $t = \frac{4\pi}{3}$ y $t = \frac{5\pi}{3}$.

$\cos t = 0$: $t = \frac\pi2$ y $t = \frac{3\pi}{2}$.

$2\sin t + 1 = 0$ significa $\sin t = -\frac12$: partiendo de
$a = -\frac\pi6$, en $\intco{0}{2\pi}$: $t = \frac{7\pi}{6}$ y
$t = \frac{11\pi}{6}$.
\end{solution}

\begin{solution}{exo:g11:trigo:8}
Por el \cref{thm:g11:trigo:equations}, $\cos 2t = \cos t$ significa
$2t = t + 2k\pi$ o $2t = -t + 2k\pi$, es decir, $t = 2k\pi$ o
$t = \frac{2k\pi}{3}$ ($k \in \Z$). La primera familia está contenida en
la segunda (tomando $k$ múltiplo de $3$), así que el conjunto de
soluciones es
\[
S = \left\{\frac{2k\pi}{3} : k \in \Z\right\}
\]
o, sobre la circunferencia, los tres puntos $M(0)$,
$M\!\left(\frac{2\pi}{3}\right)$ y $M\!\left(\frac{4\pi}{3}\right)$.
\end{solution}

\begin{solution}{exo:g11:trigo:9}
Desarrollando los dos cuadrados y usando $\cos^2 x + \sin^2 x = 1$:
\[
(\cos^2 x + 2\sin x\cos x + \sin^2 x)
+ (\cos^2 x - 2\sin x\cos x + \sin^2 x)
= 1 + 1 = 2 .
\]
\end{solution}

\begin{solution}{exo:g11:trigo:10}
Rodar sin deslizar significa que la longitud del arco es igual a la
distancia recorrida: $r\theta = 150$ cm da
$\theta = \frac{150}{30} = 5$ rad $= 5 \times \frac{180}{\pi} \approx
286.5^\circ$. Una vuelta completa hace avanzar la bicicleta la longitud de
la circunferencia, $2\pi \times 30 = 60\pi \approx 188.5$ cm, o sea, unos
$1.88$ m.
\end{solution}

\begin{solution}{exo:g11:trigo:11}
Sobre la circunferencia unidad, los puntos de \omterm{def:g10:coordgeom:system}{abscisa} exactamente
$\frac12$ son $M\!\left(\pm\frac\pi3\right)$. Los puntos de \omterm{def:g10:coordgeom:system}{abscisa}
\emph{como mucho} $\frac12$ forman el arco situado a la izquierda de la
recta vertical $x = \frac12$, que se recorre cuando $t$ va de $\frac\pi3$
hasta $\pi$, o de $-\pi$ hasta $-\frac\pi3$. Por tanto, en
$\intoc{-\pi}{\pi}$:
\[
S = \intoc{-\pi}{-\frac{\pi}{3}} \cup \intcc{\frac{\pi}{3}}{\pi}.
\]
(Los extremos $\pm\frac\pi3$ están incluidos porque la desigualdad no es
estricta.)
\end{solution}

\begin{solution}{pb:g11:trigo:1}
\textbf{1.} $30^\circ = \frac{\pi}{6}$; $135^\circ = \frac{3\pi}{4}$;
$300^\circ = \frac{5\pi}{3}$. Y $\frac{3\pi}{4} = 135^\circ$;
$\frac{5\pi}{6} = 150^\circ$.

\textbf{2.} Longitud $r\,t = 3 \times 2 = 6$ m. Los \omterm{def:g11:trigo:radian}{radianes} convierten la
longitud del arco en un producto sin más; los grados arrastrarían un
$\frac{\pi}{180}$ a cada fórmula.

\textbf{3.} \omterm{def:g11:trigo:radian}{Radianes} girados $= \frac{\text{distancia}}{\text{radio}} =
\frac{1000}{0.30} \approx 3\,333$ rad; dividiendo entre $2\pi$, unas $530$
vueltas completas.

\textbf{4.} $\cos\frac{2\pi}{3} = -\frac12$;
$\sin\left(-\frac\pi4\right) = -\frac{\sqrt2}{2}$;
$\frac{19\pi}{6} - 2\pi = \frac{7\pi}{6}$, luego
$\cos\frac{19\pi}{6} = \cos\frac{7\pi}{6} = -\frac{\sqrt3}{2}$.

\textbf{5.} $\cos^2 t = 1 - \frac{9}{25} = \frac{16}{25}$; en el segundo
cuadrante el \omterm{def:g11:trigo:cossin}{coseno} es negativo: $\cos t = -\frac45$, y
$\tan t = \frac{3/5}{-4/5} = -\frac34$.

\textbf{6.} En $t$ minutos la noria gira
$\frac{t}{30} \times 2\pi = \frac{\pi t}{15}$. En $t = 0$:
$h = 65 - 60 = 5$ m, el andén de embarque, abajo del todo (la altura del
eje menos el radio). El signo menos te coloca \emph{por debajo} del eje
cuando el ángulo girado es $0$.

\textbf{7.} $h(7.5) = 65 - 60\cos\frac{\pi}{2} = 65$ m (la altura del
eje); $h(15) = 65 + 60 = 125$ m (lo más alto);
$h(20) = 65 - 60\cos\frac{4\pi}{3} = 65 + 30 = 95$ m.

\textbf{8.} $65 - 60\cos\frac{\pi t}{15} = 95$ da
$\cos\frac{\pi t}{15} = -\frac12$: la primera solución es
$\frac{\pi t}{15} = \frac{2\pi}{3}$, es decir, $t = 10$ minutos.

\textbf{9.} $\cos\frac{\pi t}{15} \leq -\frac12$ para
$\frac{\pi t}{15} \in \intcc{\frac{2\pi}{3}}{\frac{4\pi}{3}}$, es decir,
$t \in \intcc{10}{20}$: diez minutos del viaje por encima de $95$ m,
simétricos respecto de la cima en $t = 15$.

\textbf{10.} Primer minuto:
\[
h(1) - h(0) = 60\left(1 - \cos\frac{\pi}{15}\right)
\approx 1.3 \text{ m}.
\]
Entre los minutos $7$ y $8$:
\[
60\left(\cos\frac{7\pi}{15} - \cos\frac{8\pi}{15}\right)
\approx 12.5 \text{ m},
\]
casi diez veces más. La altura cambia más deprisa al pasar por la altura
del eje y se estanca cerca del suelo y de la cima: la sinusoide es
empinada en su parte \omterm{def:g10:stats:mean}{media} y plana en sus extremos.

\textbf{11.} Radio $= \frac{40 - 4}{2} = 18$ m, eje a $4 + 18 = 22$ m y
periodo $12$ min: $h(t) = 22 - 18\cos\left(\frac{\pi t}{6}\right)$.

\textbf{12.} \omterm{def:g10:functions:extrema}{Máximo} $7 + 3 = 10$ m cuando $\sin\frac{\pi t}{6} = 1$: la
primera vez en $t = 3$ (a las 3 de la madrugada); \omterm{def:g10:functions:extrema}{mínimo} de $4$ m por
primera vez en $t = 9$; periodo $\frac{2\pi}{\pi/6} = 12$ horas.

\textbf{13.} $\sin\frac{\pi t}{6} \geq \frac12$ para
$\frac{\pi t}{6} \in \intcc{\frac\pi6}{\frac{5\pi}{6}}$:
$t \in \intcc{1}{5}$, una ventana de cuatro horas de la 1 a las 5 de la
madrugada, que vuelve a abrirse un periodo después,
$t \in \intcc{13}{17}$ (de la 1 a las 5 de la tarde).

\textbf{14.} La simetría respecto del eje vertical lleva el punto de
ángulo $t$ al de ángulo $\pi - t$ y conserva la ordenada:
$\sin(\pi - t) = \sin t$. El medio giro alrededor del centro lleva $t$ a
$\pi + t$ y cambia el signo de las dos \omterm{def:g10:coordgeom:system}{coordenadas}:
$\cos(\pi + t) = -\cos t$. Y la simetría respecto de la diagonal
intercambia \omterm{def:g10:coordgeom:system}{abscisa} y ordenada:
$\cos\left(\frac\pi2 - t\right) = \sin t$, el ``co'' de los ángulos
complementarios.

\textbf{15.} $\sin(-t) = -\sin t$: \omterm{def:g11:func:parity}{impar} (su \omterm{def:g10:functions:graph}{gráfica} tiene simetría de
medio giro); $\cos(-t) = \cos t$: par (simetría de espejo). Y
$\sin t + t^3$ es suma de dos funciones \omterm{def:g11:func:parity}{impares}: \omterm{def:g11:func:parity}{impar}.

\textbf{16.} Factorizamos: $(2\sin t + 1)(\sin t - 1) = 0$:
$\sin t = 1$ da $t = \frac\pi2$; $\sin t = -\frac12$ da
$t = \frac{7\pi}{6}$ y $t = \frac{11\pi}{6}$. Tres soluciones sobre la
circunferencia.

\textbf{17.} $1$ rad $= \frac{180}{\pi} \approx 57.3^\circ$. Así que
$\sin(1) \approx 0.841$, mientras que $\sin(1^\circ) \approx 0.0175$: un
factor de casi $50$. Moraleja: comprueba el indicador de modo antes de
fiarte de cualquier resultado trigonométrico.

\textbf{18.} La \omterm{def:g10:coordgeom:system}{abscisa} es igual a la ordenada sobre la diagonal:
$t = \frac\pi4$ y $t = \frac{5\pi}{4}$.

\textbf{19.} $\sin(0.1) = 0.09983\ldots$ frente a $0.1$: error relativo de
un $0.17\,\%$ aproximadamente. Para ángulos pequeños, la cuerda se abraza
al arco, y el límite $\frac{\sin t}{t} \to 1$ del año que viene es
exactamente esta observación convertida en teorema.

\textbf{20.} \omterm{def:g11:trigo:radian}{Radianes}: el ángulo es el arco sobre la circunferencia
unidad, así que los ángulos se calculan como longitudes. \omterm{def:g11:trigo:cossin}{Coseno} y \omterm{def:g11:trigo:cossin}{seno}:
las \omterm{def:g10:coordgeom:system}{coordenadas} de la aguja del reloj, la sombra en el suelo y la altura
en la pared. Simetrías de la circunferencia: todo el catálogo de
identidades, leído en espejos y medios giros. Ecuaciones $\cos t = c$: el
horario del reloj, con la circunferencia mostrando todas las soluciones a
la vez. Disfraces: la noria llevaba $65 - 60\cos$ y la marea, $7 + 3\sin$;
el mismo reloj, las mismas matemáticas.
\end{solution}
